\section{A finite-difference benchmark for hyperbolic heat conduction}\label{sec:schemes}

Hyperbolic heat laws are usually validated in the time domain against a smooth test problem. That choice hides the real difficulty: a front that the scheme has to carry without spurious oscillation or excessive smearing. A small benchmark with an exact solution is given, and what three standard schemes achieve is reported.

\subsection{Problem}
The uncoupled Cattaneo system in first-order form,
\begin{equation}
\rho c\,\theta_t+q_x=0,\qquad \tau\,q_t+q=-k\,\theta_x ,
\end{equation}
on $x>0$ with $\theta=q=0$ initially and $q(0,t)=q_0H(t)$, is solved in units $\rho c=k=\tau=q_0=1$, so $c_T=1$ and the front is at $x=t$. The exact solution is the inverse Laplace transform of \cref{eq:pulse}, taken as zero ahead of the front. At $t=5$ it gives $\theta=1.7615,\ 0.5923,\ 0.0959$ at $x=1,\,3,\,4.9$ and the jump at the front is $q_0e^{-x/2\ell}/(\rho cc_T)=e^{-2.5}=0.082$.

\subsection{Schemes}
The system is advanced by Strang splitting. The source $-q/\tau$ is integrated exactly over half steps, and the hyperbolic part (characteristic speeds $\pm c_T$) by one of three schemes: (i) Lax--Friedrichs, (ii) first-order upwind on the characteristic variables, (iii) second-order Lax--Wendroff. All three run at Courant number $0.9$, and the boundary value $q=q_0$ enters through a ghost cell. The code is the finite-difference benchmark computation listed in \cref{app:repro}.

\subsection{Results}
Errors are $L^1$ norms over the sampled interval at $t=5$. Over $0.05\le x\le4.95$, which contains the front, the observed orders between $N=800$ and $1600$ cells are $0.85$ for Lax--Friedrichs, $0.94$ for upwind and $0.58$ for Lax--Wendroff, with errors at $N=1600$ of $2.7\times10^{-3}$, $1.3\times10^{-3}$ and $8.6\times10^{-4}$. Away from the front ($0.05\le x\le3.5$) the picture changes:
\begin{center}
\begin{tabular}{@{}lccccc@{}}\toprule
$N$ & 100 & 200 & 400 & 800 & 1600\\\midrule
upwind $L^1$ & $9.6\times10^{-3}$ & $4.3\times10^{-3}$ & $2.1\times10^{-3}$ & $1.1\times10^{-3}$ & $5.3\times10^{-4}$\\
Lax--Wendroff $L^1$ & $1.1\times10^{-3}$ & $2.7\times10^{-4}$ & $6.1\times10^{-5}$ & $1.5\times10^{-5}$ & $3.8\times10^{-6}$\\
\bottomrule\end{tabular}
\end{center}
with observed orders $1.0$ for the upwind scheme and $2.0$ for Lax--Wendroff in the smooth region. Formal accuracy is reached away from the front, and the neighborhood of the discontinuity controls the global error, where all three schemes converge at an order below one. A single global error norm therefore says little about how a scheme carries a front. A benchmark for hyperbolic heat conduction should report the error around the front separately from the error away from it.

\begin{figure}[t!]\centering\small
\input{figs_tikz/fig10_schemes.tex}
\caption{Original schematic; the setting follows \citep{Joseph1990}: finite-difference benchmark for hyperbolic heat conduction -- the front at $x=t$ must be carried without smearing or oscillation; error measured in a window around and away from the front.}\label{fig:schemes-geo}
\end{figure}

A domain of influence theorem for generalized thermoelasticity with memory-dependent derivative was proved by \citep{Sarkar2019a}.

\subsection{A caution about the reference solution}
The exact solution above was obtained by numerical Laplace inversion, and the choice of inversion algorithm matters at the accuracy required here. At $x=1$, $t=5$ the Talbot method with default settings gave $1.76126$ at 30 working digits and $1.76086$ at 20 digits, whereas the de Hoog and Stehfest methods give $1.761515$ and $1.761516$ at 20 digits. The finite-difference solution at $N=800$ agrees with the latter to $2\times10^{-6}$. A reference solution with an error of $3\times10^{-4}$ or more would have hidden the second-order convergence of Lax--Wendroff at $N\gtrsim400$. Inversion results used as benchmarks should be cross-checked with a second algorithm (\cref{sec:inversion}).

\begin{summarybox}[Summary]
The Cattaneo benchmark has an exact front at $x=t$ (in units $c_T=1$). Lax--Friedrichs, upwind and Lax--Wendroff all converge below order one globally because of the discontinuity; away from the front the orders are $1.0$ and $2.0$ respectively. Report front-window and smooth-window errors separately, and cross-check Laplace-inversion reference solutions with a second algorithm.
\end{summarybox}
